% Laget automatisk av elevark/lag_elevark.py fra sim/ekvivalensprinsippet.html. Endringer her blir overskrevet.
\documentclass[11pt]{article}
\usepackage{elevark}
\begin{document}
\elevarkhode{Fysikk 2 · Generell relativitetsteori}{Ekvivalensprinsippet}{https://fysikklab.no/sim/ekvivalensprinsippet.html}{Simuleringen}

Du våkner i en lukket heis uten vinduer og kjenner at du står på gulvet. Står heisen på jorda, eller er den i en rakett som akselererer ute i rommet? Einstein mente at ingen forsøk inne i heisen kan avgjøre det. Det kalles ekvivalensprinsippet, og det fører til at lys må bøyes av tyngdekraften.

\section*{Oppgaver}

\begin{oppgaveliste}
\oppgave \fase{gjett}{Gjett}En ball kastes vannrett i en rakett som akselererer med 9,81~m/s² langt ute i rommet. Følger ballen en rett linje eller en krum bane, sett fra en astronaut inne i raketten?
\svar
\oppgave \fase{test}{Test}Sammenlign «I ro på jorda» og «Akselererer i rommet». Er det noen forskjell på banene sett innenfra? Hva er forskjellen sett utenfra?
\svar
\oppgave \fase{forklar}{Forklar}Sammenlign «Faller fritt» og «Svever i rommet». Hvorfor kan astronautene i romstasjonen sveve, selv om tyngdekraften der oppe er nesten like sterk som på bakken?
\svar
\oppgave \fase{forklar}{Forklar}Se på lyspulsen i raketten, sett utenfra. Hvorfor bøyer lyset seg, sett innenfra? Bruk ekvivalensprinsippet til å forklare hvorfor lys også må bøyes av tyngdekraften på jorda.
\svar
\oppgave \fase{test}{Regn}Regn ut hvor mye lyset faller på vei over heisen på 3,0~m med ekte lysfart. Hvorfor merker vi aldri dette i hverdagen?
\svar
\oppgave \fase{utfordring}{Utfordring}Ekvivalensprinsippet sier at «svevende» og «uten tyngdekraft» ikke kan skilles fra hverandre lokalt. Hvorfor gjelder det bare i et lite område? Tenk på en svært stor heis som faller mot jorda.
\svar
\end{oppgaveliste}

\end{document}
